\section{Introduction}
\label{sec:intro}

Biological collections are the material foundation of taxonomy and of one of
the great tasks of the life sciences---documenting the Earth's biodiversity.
According to Index Herbariorum, the world holds roughly 3\,500 active herbaria
containing $\sim$398 million physical specimens~\cite{thiers2024ih}. These
collections underpin systematic, ecological, climatic, and conservation
research and serve as a verifiable record of the presence of taxa in space and
time~\cite{davis2023herbarium}. The Komarov Botanical Institute (LE) holds the
largest herbarium in Russia and the fifth largest worldwide, with more than
7.3 million specimens accumulated over 300 years; its mycological fund
($\sim$350\,000 specimens) is the largest fungal collection in the country and
includes numerous type specimens of Ustilaginales.

\paragraph{The dark-data problem.}
Despite their scientific value, collection metadata remain ``dark data.'' At
LE, only $\sim$5\,\% of vascular plants are imaged, the figure for the
mycological part is below 1\,\%, and structured metadata exist for only
$\sim$22\,\% of specimens. Globally, the rate of metadata transcription has not
increased in over a decade and remains the principal bottleneck of
biodiversity mobilization into open databases~\cite{owen2020nlp}. Manual
transcription is expensive, slow, and noisy due to the heterogeneity of
curatorial schools and epochs. The resulting data should follow the FAIR
principles and biodiversity standards such as Darwin
Core~\cite{wieczorek2012dwc}.

\paragraph{Why mycological type labels are hard.}
Our working corpus is the smut-fungi (Ustilaginomycotina, ``Smuts'')
subcollection of LE: 394 type-specimen images with inventory prefixes
\texttt{LE~F-*}. A first visual analysis reveals extreme heterogeneity along
several axes (Tab.~\ref{tab:corpus}): epochs from 1873 to 2009; four languages
(Latin scientific names, German, English, and Russian including pre-reform
orthography); four text types (typeset, typewritten, handwritten, mixed);
diverse date formats; legacy toponyms (e.g.\ ``USSR: Kazachstaniae, Alatau
Transiliense'') alongside modern DMS coordinates; visually expressed status
markers (\emph{Isotypus}, \emph{Holotypus} colored stickers); and complex
layouts (multi-label sheets, barcodes, envelope spreads). Crucially, none of
the closest prior systems is validated on material of such temporal and
linguistic breadth, and most target \emph{plant} collections that rely on the
World Flora Online rather than fungal
nomenclature~\cite{turnbull2025hespi,wen2025ocrllm}.

\begin{table}[t]
\centering
\small
\setlength{\tabcolsep}{4pt}
\begin{tabular}{@{}ll@{}}
\toprule
Axis & Range / examples \\
\midrule
Epochs & 1873--2009 (Sydow 1898; Vánky 1988--96) \\
Languages & Latin, German, English, Russian (pre-reform) \\
Text type & typeset, typewritten, handwritten, mixed \\
Dates & \texttt{20.\,10.\,1898}, \texttt{13.\,VI.\,1946}, \texttt{26.II.1988} \\
Geography & legacy toponyms; DMS coordinates \\
Status & Isotypus / Holotypus (colored stickers) \\
Layout & multi-label sheets, barcodes, envelope spreads \\
\bottomrule
\end{tabular}
\caption{Heterogeneity of the smut-fungi (Smuts) subcollection. This diversity
is the scientific motivation of our study: it stresses every stage of the
pipeline, from orientation to handwriting recognition to taxonomic
normalization.}
\label{tab:corpus}
\end{table}

\paragraph{Goal.}
We aim to build and empirically validate a modular pipeline that extracts
structured, machine-readable metadata from mycological label images, aligned
with a 21-field Darwin~Core-compatible schema~\cite{wieczorek2012dwc} extended
with mycology-specific fields (host plant, type status), with per-field
confidence and an audit trail suitable for curator review. Our study is
deliberately incremental: each step is compared against the previous one on the
same gold set, so that the contribution of every component can be attributed
quantitatively.

\paragraph{Contributions.}
\begin{enumerate}[itemsep=1pt,topsep=2pt,leftmargin=1.4em]
\item A rigorous, reproducible benchmark of \textbf{17 experiments} on a real
mycological type corpus, from a seven-engine classical OCR ensemble to modern
VLMs, all evaluated with a single unified metric suite
(\cref{sec:experiments}).
\item Evidence that \textbf{domain adaptation dominates architecture}: a LoRA
adapter (rank~32, $0.5\,\%$ trainable parameters) trained on 310 examples
raises $F_1$ from $0.807$ (zero-shot VLM) to \textbf{0.972} on the full corpus,
outscoring 18 stronger or larger models (\cref{sec:results}).
\item A detailed \textbf{LoRA hyperparameter study} (83 configurations)
quantifying the effect of learning rate, rank, $\alpha$, epochs, and target
modules, with a compact recipe.
\item A practical \textbf{production recipe}: adapter merging plus batched
inference gives a $5\times$ speedup ($2.77$\,s/image) with no measurable quality
loss; we report throughput for corpus-scale deployment.
\item Two carefully designed \textbf{negative results}---a taxonomic-API
validator is a signal, not an $F_1$ booster, and a naive LLM corrector degrades
quality---plus a \textbf{cross-model GT audit} that finds 51 likely annotation
omissions.
\item A \textbf{generalization test} on 534 unseen \emph{Anthracoidea}
specimens, confirming a $99.8\,\%$ parse rate and consistent per-field behavior.
\end{enumerate}
